\documentclass[14pt,a4paper,oneside]{article}

\usepackage[utf8]{inputenc}   


\usepackage{amsmath,amsthm,amssymb,appendix,bm,mathrsfs,ctex,geometry,tcolorbox,fancyhdr,tikz,esint,appendix}


%% GB/T 7714 宏包 ---- 引文格式
\usepackage[sort&compress]{gbt7714}
\bibliographystyle{gbt7714-numerical}



\numberwithin{equation}{subsection}             %%%%%% 公式编号




\geometry{left=1.7cm, right=1.7cm, top=2.54cm, bottom=2.54cm}    %%%%% 页面尺寸
\linespread{1.5}                             %%%%% 行距1.5倍
\pagestyle{fancy}



%%    设置页眉页脚
\fancyhf{}
\renewcommand{\headrulewidth}{0.15mm}
\fancyhead[C]{文献阅读报告}                   %%%%%% 页眉
\fancyfoot[C]{\thepage}                      %%%%%% 页脚 页码
\setlength{\headheight}{15pt}


%%   设置各种链接颜色为黑色
\usepackage[colorlinks,
            linkcolor=black,
            anchorcolor=black,
            citecolor=black,
            urlcolor=black
            ]{hyperref}



%%   例子编号
\newcounter{example}
\counterwithin{example}{subsection}
\newcommand{\examplename}{Example}
\newcommand{\exampleTitle}[1]{%
\textbf{\refstepcounter{example}\examplename\theexample:#1}\par}


%%  插入图片配置
\usepackage{graphicx,subfigure}
\usepackage[]{caption2} %新增调用的宏包
\numberwithin{figure}{subsection}               %%%%%% 图片编号

\graphicspath{ {./Figure/} }            %%%%% 设置图片文件所在目录
\DeclareGraphicsExtensions{.pdf,.svg,.jpeg,.png,.jpg}     %%%%% 支持图片格式


%%    支持tikz作图
\usepackage{tikz}
\usetikzlibrary{backgrounds,arrows}


\usepackage{pgfplots}
\pgfplotsset{compat=1.15}
\usepackage{mathrsfs}


\title{\textbf{文献阅读报告 10\\Ultrafast serrodyne optical frequency translator}}                           %%%%% 标题
\author{骆禹天}                                           %%%%% 作者
\date{2024/**/**}                                        %%%%%% 日期

%%   定义一些数学格式
\newtheorem{theorem}{定理}[subsection]
\newtheorem{definition}[theorem]{定义}
\newtheorem{lemma}[theorem]{引理}
\newtheorem{corollary}[theorem]{推论}
\newtheorem{example}[theorem]{例}
\newtheorem{proposition}[theorem]{命题}
\renewcommand{\abstractname}{\Large\textbf{摘要}}




%------------------------- 开始文档  -------------------------%



\begin{document}

\pagenumbering{Roman}         %%%%% 目录前页码统计按罗马数字

\maketitle

\setcounter{page}{1}          %%%%% 标题开始统计页数

%------------------------- 目录  ------------------------------%

\tableofcontents              %%%%% 插入目录

\newpage
\setcounter{page}{1}          %%%%% 目录后重新统计页数


\pagenumbering{arabic}        %%%%% 目录后统计页码为阿拉伯数字



%%%%%%%%%%%%%%%%%%%% 开始第一节  %%%%%%%%%%%%%%%%%%%%%


\section{Introduction}

回顾:

超快高功率激光器可提供飞秒范围的超短脉冲。随着啁啾脉冲放大 (CPA) 的发明，过去三十年在超快科学和计量学以及强场物理和材料科学方面取得了重大进展。这一进步的一个重要部分可归因于超快高峰值或平均功率激光技术的进步。应用:

· 捕获亚飞秒电子动力学，

· 深入了解质子结构，

· 短波长纳米光刻开发用于下一代芯片制造

问题: 许多创新和应用需要专用的激光参数，包括波长、强度、脉冲长度、光束特性等, 但波长是一个非常重要但不太灵活的参数。



如今通常使用镱基, 钛宝石基和最近铥基激光平台产生高峰值, 高平均功率的飞秒脉冲。这些激光系统受限于激光增益介质带宽只能在特定波长下工作. 

目前有多种选择实现波长调谐, 通常在低功率下使用:

· 光纤中的受激拉曼散射、

· 孤子频移6,7、

· 色散波生成8,9

· 空心光纤或毛细管中拉曼频移方案10,11。

· 光参量放大器中参量频率转换。

缺陷: 转换效率低或功率处理能力有限。

如何提高效率? 后压缩技术: 可高效实现几飞秒超短脉冲, 特别是与镱基激光器结合使用，可获得平均功率达千瓦的超短脉冲13，接近太瓦峰值功率范围14。

不足: 缺乏波长调谐选项。 

解决方法: Serrodyne 移频方法

历史上, 由射频技术领域使用15,16，后应用于连续激光17-21——可解决这一问题. 从连续激光到超短脉冲，该方法都可提供高效率以及同时具备太瓦级峰值功率与千瓦级平均功率，同时支持精密计量应用所需的关键特性，包括相位相干性和载波包络偏移频率维持。

锯齿差原理指出，对于给定信号，当在时域中应用线性相位时，信号的频率会发生变化22。例如，可以通过电光调制来应用这种相位。相反，我们使用利用克尔效应的全光学方法，该方法可用于将线性幅度调制及时转换为线性相位斜坡。中心频率为 ω0 的时间锯齿脉冲经过自相位调制 (SPM)，频移为 ω = ω0 + Δω。频移的大小和方向由下式给出


%--------------------  插入参考文献  --------------------%

\newpage

\clearpage           %%%  “目录”中显示的页码正确, 需要添加 \clearpage
\phantomsection      %%   “目录”中的链接能正确跳转, 需要添加 \phantomsection
\addcontentsline{toc}{section}{参考文献}   %向目录中添加条目, 以章的名义

\bibliography{ref}

\end{document}